\chapter*{Resumen}
\addcontentsline{toc}{chapter}{Resumen}

Los enfoques convencionales basados en datos para la planeación del transporte
público con frecuencia predicen dónde el transporte \emph{ya existe} en lugar de
dónde se \emph{necesita}, produciendo modelos circulares que premian el statu quo.
Esta tesis desarrolla un marco transferible y orientado a la demanda que, en
cambio, diagnostica la demanda de transporte \emph{no atendida} y genera corredores
candidatos para atenderla, aplicado a la Zona Metropolitana de Guadalajara
(\zmg; \num{1881} áreas geoestadísticas básicas, \ageb, en diez municipios).

El marco se desarrolla en ocho etapas encadenadas: (i)~una superficie de demanda de
transporte estimada a partir de datos censales, económicos y de la red vial
mediante un modelo gravitacional doblemente restringido; (ii)~la calibración del
parámetro de decaimiento por distancia contra la Encuesta Origen--Destino 2022;
(iii)~una capa de oferta que mide la accesibilidad basada en \textsc{gtfs} y un
índice independiente de \emph{brecha de cobertura} que se convierte en la variable
objetivo supervisada, eliminando la circularidad de las etiquetas de ``tiene
parada''; (iv)~un mapa de priorización Nodo--Lugar--Personas con un término
explícito de equidad; (v)~una función formal de evaluación multiobjetivo; (vi)~un
generador de corredores basado en anclas de frontera, un tronco de diámetro sobre
el árbol de expansión mínima y una compuerta de factibilidad por rectitud;
(vii)~una auditoría de la red existente; y (viii)~un conjunto de validación que
combina retro-pruebas de exclusión, comparación y métricas de cobertura antes y
después.

La capa de diagnóstico es la contribución sólida: el índice de brecha de cobertura
se corrobora fuera de muestra con la apertura en 2025 de la Línea~4 de \zmg, cuyo
corredor presenta \num{1.6} veces la tasa metropolitana de brecha alta. La capa
generativa produce al menos un corredor sustantivo, factible y meritorio, a la vez
que caracteriza honestamente sus limitaciones. Finalmente, todo el marco se
transfiere de principio a fin a dos zonas metropolitanas de tamaño contrastante
---Toluca y Aguascalientes---, demostrando que es genuinamente portable y no un
desarrollo a la medida de \zmg.

\vspace{1em}
\noindent\textbf{Palabras clave:} transporte público; demanda de viajes; brecha de
cobertura; modelo Nodo--Lugar; optimización multiobjetivo; equidad; \textsc{gtfs};
transferibilidad; Guadalajara.

\clearpage
